\section{The bounded-parent survival gap}
\label{pup:section}

By Theorem~\ref{cm:all-moments}, for every $\alpha>\lambda$ and every fixed $r>1$,
\begin{equation}\label{pup:controlled-second}
 \sup_{q\in\mathcal Q}\mathbb E_q C(t)^r\le K_{r,\alpha}e^{r\alpha t}.
\end{equation}
The constants are uniform in the newborn exterior. This estimate does not use $\lambda=\mu$.

\subsection{Uniform control while the actual parent remains bounded}

Fix $M\ge2$ and, for a newborn tag whose parent starts at degree
$d\le M$, let $\tau_M=\inf\{t:D_p(t)=M+1\}$.  Define
\[
 L_M(t)=\sup_{x:\,D_p(0)\le M}
       \mathbb E_x[C_{\rm tag}(t)\mathbf1_{\{t<\tau_M\}}].
\]
The supremum is over the actual finite exteriors, with the same permitted
inert stubs as in the earlier profile construction.

\begin{lemma}[A stopped inverse-likelihood moment]\label{pup:stopped-I}
Under the reference law replacing parent degree gains by
$h(t)=(1+t)^{-2}$, stopped at $\tau_M$, put
\[
 I_M=\int_0^{\tau_M}\lambda_p(t)\,dt,
 \qquad B_M=3M^2(1+3M).
\]
Then, for $aB_M<1$,
\begin{equation}\label{pup:stopped-moment}
 \mathbb E e^{aI_M}\le(1-aB_M)^{-1}.
\end{equation}
The expected remaining occupation is at most $B_M$ conditionally after
every stopping time before $\tau_M$.
\end{lemma}
\begin{proof}
Before stopping there are at most $M$ neighbour labels of $p$ in total.
Each real neighbour, from its time of birth, receives immigration at
rate at least $1/M$, by the permanent-slot construction.  The reciprocal
population bound gives a conditional expected remaining contribution
of at most
\[
 \int_0^\infty(1+3M)e^{-t/(3M)}dt=3M(1+3M)
\]
to the integral of its reciprocal degree.  The estimate is valid
also for a neighbour initially of degree one and for stopped integrals.
Summing over at most $M$ present or future neighbour labels gives $B_M$.
The same estimate restarts conditionally from any stopping time; the
number of remaining possible neighbour labels is still at most $M$.
Iterated ordered-time integration gives
$\mathbb E I_M^r\le r!B_M^r$, and summation proves
\eqref{pup:stopped-moment}.
\end{proof}

\begin{theorem}[The bounded-parent survival gap]\label{pup:low-survival}
If $\mu<\lambda$, then for every fixed integer $M\ge2$ there is
$\gamma_M<\lambda$ such that
\begin{equation}\label{pup:low-survival-bound}
 \limsup_{t\to\infty}\frac1t\log L_M(t)\le\gamma_M.
\end{equation}
Thus an exponent gap cannot be supported by newborn histories whose
parent remains below a fixed degree threshold throughout observation.
\end{theorem}
\begin{proof}
Fix $\varepsilon\in(0,1)$ and split the event $t<\tau_M$ according
to whether a parent degree gain occurs after time $\varepsilon t$.
Theorem~\ref{ps:stabilization} gives
\begin{gather*}
 \Pr\{\text{such a late gain before or at }\tau_M\}
 \le Q_M(\varepsilon t),\\
 Q_M(r)=O_M(\operatorname{poly}(r)e^{-\kappa_M r}),
 \qquad \kappa_M=1/(3M).
\end{gather*}
By \eqref{pup:controlled-second}, the contribution of these histories
has upper exponent at most $\alpha-\kappa_M\varepsilon/2$ for every
$\alpha>\lambda$.

Let $A_t$ be the other event: $t<\tau_M$ and all parent degree gains
through time $t$ occur by $\varepsilon t$.  Theorem~\ref{plc:theorem} bounds the actual expectation on $A_t$ by a fixed polynomial
in $t$ times the reference expectation, uniformly in the initial
exterior.  Only the finitely many starting degrees $d=2,\ldots,M$
need be considered.

Realize each reference marginal by modifying a fixed genuine finite
star completion of the newborn tag.  Let $Z_t$ be its actual-to-reference
likelihood, and let $N(r)$ be its full reference real population, stopped
at $\tau_M$.  Before stopping,
\[
 \lambda_p(r)\ge1/D_u(r)\ge1/N(r),\qquad h(r)\le1.
\]
There are at most $M$ parent degree gains.  Hence on $A_t$,
\[
 Z_t^{-1}
 =\exp\!\left(\int_0^t(\lambda_p-h)dr\right)
       \prod_{\sigma_j\le t}\frac{h(\sigma_j)}{\lambda_p(\sigma_j-)}
 \le e^{I_M}N(\varepsilon t)^M.
\]
Choose $a>0$ with $2aB_M<1$.  By Cauchy--Schwarz and
Lemma~\ref{pup:stopped-I},
\begin{equation}\label{pup:prefix-inverse}
 \mathbb E_{\rm ref}[Z_t^{-a}\mathbf1_{A_t}]
 \le K_{a,M}e^{aM\varepsilon t}.
\end{equation}
Indeed, the reference total birth rate is at most $N+1$; a shifted Yule
domination gives $\mathbb E N(r)^v\le K_{M,v}e^{vr}$ for every fixed
$v>0$.

Put $\theta=a/(a+2)$, so $a=2\theta/(1-\theta)$.  Three-factor
H\"older now gives
\begin{align*}
 \mathbb E_{\rm ref}[C(t)\mathbf1_{A_t}]
 &\le
 \bigl(\mathbb E_{\rm ref}[Z_tC(t)\mathbf1_{A_t}]\bigr)^\theta
 \bigl(\mathbb E_{\rm ref}C(t)^2\bigr)^{(1-\theta)/2}\\
 &\qquad\times
 \bigl(\mathbb E_{\rm ref}[Z_t^{-a}\mathbf1_{A_t}]\bigr)^{(1-\theta)/2}.
\end{align*}
The first expectation is at most the total core mean from the fixed
genuine star completion, and therefore has upper exponent $\mu$.
The second has upper exponent $2\alpha$ by the controlled-moment
bound; stopping only decreases its core count.  Formula
\eqref{pup:prefix-inverse} bounds the third.  After the polynomial
forward likelihood factor, the good contribution consequently has
upper exponent
\[
 \theta\mu+(1-\theta)\alpha+\theta M\varepsilon.
\]
Take $a=1/(4B_M)$, so $\theta=1/(8B_M+1)$, and write
$g=\lambda-\mu>0$.  Choose $\varepsilon=g/(4M)$ and then
$\alpha=\lambda+\delta$ with
\[
 0<\delta<
 \min\left\{\frac{\theta g}{4(1-\theta)},
              \frac{\kappa_M\varepsilon}{4}\right\}.
\]
Both displayed contribution exponents are strictly less than $\lambda$:
the good one is below $\lambda-\theta g/2$, and the bad one below
$\lambda-\kappa_M\varepsilon/4$.  The finite maximum over $d$ changes
only constants.  This proves \eqref{pup:low-survival-bound}.
\end{proof}

\begin{corollary}[A bounded retained root is also subcritical]
\label{pup:root-stay}
Under $\mu<\lambda$, for every fixed integer $H\ge2$,
\[
 \limsup_{t\to\infty}\frac1t\log\sup_x
 \mathbb E_x[C_{\rm tag}(t)\mathbf1_{\{D_u(t)\le H\}}]<\lambda.
\]
\end{corollary}
\begin{proof}
Partition the tagged core genealogy into the root and the descendant
components of its future children. On $D_u(t)\le H$, each child's
parent has stayed at degree at most $H$ since its birth, because degree
is nondecreasing. Conditional on the child's birth, its weighted
remaining mean is therefore at most $L_H$ at the remaining horizon.
The root produces children at intensity at most one. The compensator
formula and positivity give
\[
 \sup_x\mathbb E_x[C_{\rm tag}(t)\mathbf1_{\{D_u(t)\le H\}}]
 \le1+\int_0^tL_H(v)\,dv.
\]
Theorem~\ref{pup:low-survival} bounds the right side with an exponent
strictly below $\lambda$.
\end{proof}

\subsection{Removing arrivals after a high-degree crossing}

For an initially low-degree parent, define the hybrid process to agree
with the actual process up to $\tau_M$, and thereafter delete all
parent-to-tag arrivals.  Outward tag selections and all other exterior
transitions may continue; the retained component after $\tau_M$ follows
the inert law.  For an initially high real parent, delete its incoming
channel from time zero.  For an initially inert parent stub, set
$\tau_M=0$ and use inert evolution from the start.  Let $B_M^0(t)$ be the supremum of the hybrid
retained means over all newborn environments.

\begin{theorem}[A subcritical crossing truncation]\label{pup:hybrid}
Suppose $g=\lambda-\mu>0$.  If $M\ge2$ is fixed and
$\delta=1/(M+1)<g$, then
\begin{equation}\label{pup:hybrid-gap}
 \limsup_{t\to\infty}\frac1t\log B_M^0(t)<\lambda.
\end{equation}
\end{theorem}
\begin{proof}
An initially high or inert parent gives the inert newborn mean, with exponent
$\mu$.  For an initially low parent, histories that have not crossed by
$t$ are covered by Theorem~\ref{pup:low-survival}.  Choose
$\gamma<\lambda$ and $K<\infty$ such that $L_M(r)\le Ke^{\gamma r}$.

Split the crossing histories at $\tau_M=\varepsilon t$, where
$\varepsilon>0$ will be small.  The expected core count carried by
crossings in $dr$ is at most $M L_M(r)\,dr$: the pre-crossing parent
degree-gain intensity is at most $M$, and the crossing leaf is added
outside the retained component.  After the crossing the conditional
mean is at most $C(\tau_M)H(t-\tau_M)$.  Thus the late-crossing part
is at most
\[
 M\int_{\varepsilon t}^t L_M(r)H(t-r)\,dr.
\]
For every $\alpha>\lambda$, $H(v)\le K_\alpha e^{\alpha v}$.
The last display has upper exponent at most
$\alpha-(\alpha-\gamma)\varepsilon$, which is below $\lambda$ when
$\alpha$ is sufficiently close to $\lambda$.

For early crossings, use the same finite reference completion and
polynomial forward likelihood reduction as before.  Change parent
degree gains to $h(t)$ only until $\tau_M$; thereafter restore all
genuine transition rates except the deleted parent-to-tag channel.
The retained-and-crossing-time marginal is autonomous: it is the
previous reference law before crossing, then inert evolution.
Its full reference population has birth rate at most $N+1$.

Relative to the fixed genuine finite-tree law, the reference-path
likelihood includes an additional survival factor for the deleted
channel after crossing.  Since the parent's degree is then at least
$M+1$, its missing intensity is at most $\delta$.  On
$\tau_M\le\varepsilon t$, the inverse likelihood therefore obeys
\[
 Z_t^{-1}\le e^{I_M+\delta t}N(\varepsilon t)^M.
\]
With the same $a>0$ and $\theta=a/(a+2)$ as in the preceding proof,
\[
 \mathbb E_{\rm ref}
 [Z_t^{-a};\tau_M\le\varepsilon t]
 \le K e^{a(\delta+M\varepsilon)t}.
\]
Here the density is for the actual law restricted to having no arrivals
in the deleted channel after crossing; its weighted core expectation
is consequently bounded by the unrestricted genuine finite-tree core
mean.  The same three-factor H\"older argument shows that the early
crossing contribution has upper exponent at most
\[
 \theta\mu+(1-\theta)\alpha+
       \theta(\delta+M\varepsilon).
\]
Choose $\varepsilon>0$ with $\delta+M\varepsilon<g$, and then
$\alpha>\lambda$ sufficiently close to $\lambda$ for this exponent
and the late-crossing exponent both to be below $\lambda$.
All forward likelihood factors are polynomial in $t$ and there are
only finitely many starting parent degrees.  This proves the claim.
\end{proof}

